\chapter{Ask-Jade Pipeline}

% ═══════════════════════════════════════════════════════════════════════════════
\section{Introduction}
% ═══════════════════════════════════════════════════════════════════════════════

This chapter describes the technical architecture of the Ask-Jade system.
It presents the ingestion pipeline, which transforms raw documents into
a searchable knowledge base, and the retrieval pipeline, which processes
user queries and generates grounded, cited responses. The rationale behind
each technological choice is also discussed.

% ═══════════════════════════════════════════════════════════════════════════════
\section{Pipeline Architecture Overview}
% ═══════════════════════════════════════════════════════════════════════════════

To address the challenges described in previous chapters, a modular
Retrieval-Augmented Generation (RAG) pipeline was designed to structure
and exploit unstructured documents. LlamaIndex~\cite{llamaindex} is used
as the main orchestration framework, connecting all stages from ingestion
to answer generation. Based on the comparison in Chapter 2, it was chosen
over LangChain due to its better support for indexing, vector storage,
and query management.

The pipeline is organised into two complementary layers:

\begin{itemize}
  \item The \textbf{Ingestion pipeline} operates offline, before any user
        interaction. It collects documents, parses and normalises their
        content, splits text into semantic chunks, generates embeddings, and
        persists everything in the vector database.

  \item The \textbf{Retrieval pipeline} operates at query time via the
        LlamaIndex Condense Plus Context Chat Engine. It rewrites
        the user query from conversation history, expands it into multiple
        variants, retrieves relevant chunks, merges results via RRF Fusion,
        reranks candidates, and produces a grounded, cited answer using a
        configurable Large Language Model. Every generation trace is saved to
        Langfuse for observability.
\end{itemize}

%Table~\ref{tab:pipeline_stages} summarises each stage with its input, output,
%and responsible component.

%\begin{table}[H]
%\centering
%\caption{Stages of the Ask-Jade Pipeline}
%\label{tab:pipeline_stages}
%\renewcommand{\arraystretch}{1.4}
%\begin{tabular}{|l|l|l|l|}
%\hline
%\textbf{Layer} & \textbf{Stage} & \textbf{Input} & \textbf{Output / Component} \\
%\hline
%\multirow{4}{*}{Ingestion}
%  & Document acquisition   & Raw files             & Parsed text \textemdash{} Docling \\
  %& Chunking               & Parsed text           & Text chunks \textemdash{} LlamaIndex \\
 % & Embedding              & Chunks                & Dense vectors \textemdash{} BGE-M3 \\
%  & Storage                & Vectors + metadata    & PgVector \\
%\hline
%\multirow{6}{*}{Retrieval}
  %& Query rewriting        & Query + history       & Reformulated query (Condense Step) \\
 % & Multi-query generation & Reformulated query    & 2--3 query variants \\
 % & Semantic retrieval     & Query variants        & Top-10 candidates per variant \\
 % & RRF Fusion             & Top-10 $\times$ N lists & Merged \& ranked candidates \\
  %& Reranking              & Merged candidates     & Top-5 chunks \textemdash{} BGE-reranker-v2-m3 \\
 % & Answer generation      & Top-5 + query         & Cited answer \textemdash{} LLM (configurable) \\
%\hline
%\end{tabular}
%\end{table}

The following sections detail each pipeline in turn.

% ═══════════════════════════════════════════════════════════════════════════════
\section{Ingestion Pipeline}
\label{sec:ingestion}
% ═══════════════════════════════════════════════════════════════════════════════

The ingestion pipeline transforms raw, heterogeneous documents into structured,
searchable knowledge. It covers four stages: acquisition, parsing and chunking,
embedding generation, and vector storage. Figure~\ref{fig:ingestion_pipeline}
illustrates this flow.

%\begin{figure}[H]
  %\centering
  %\includegraphics[width=\textwidth]{figure_ingestion.jpg}
  %\caption{Ingestion pipeline of the Ask-Jade system}
  %\label{fig:ingestion_pipeline}
%\end{figure}

% ─────────────────────────────────────────────────────────────────────────────
\subsection{Document Acquisition}
% ─────────────────────────────────────────────────────────────────────────────
 
The system currently supports four document formats:
PDF, DOCX and PPTX.
All documents come from the SharePoint connector,
which automatically gets files from our organisations repository
using a special LlamaIndex connector. This way new and updated corporate documents
are always added automatically without needing any help.
 
Prior to entering the pipeline, each document is checked for near-duplicates
via TLSH (Trend Micro Locality Sensitive Hash). Unlike cryptographic
hashes such as SHA-256, which only detect exact byte-for-byte duplicates,
TLSH computes a fuzzy fingerprint over the document content and measures
similarity via a distance score. Two documents whose TLSH distance falls below
a defined threshold are considered near-duplicates; the incoming document is
then discarded. Otherwise it is forwarded to the next stage. This approach
prevents redundant storage even when documents differ only by minor edits,
metadata updates, or formatting changes.

% ─────────────────────────────────────────────────────────────────────────────
\subsection{Document Parsing and Chunking}
% ─────────────────────────────────────────────────────────────────────────────

DOCX files are first converted to PDF to standardise the processing workflow.
All documents are then processed by Docling~\cite{docling},
an open-source document-understanding library developed by IBM Research.
Table \textbf{4.1} compares Docling with
PyMuPDF~\cite{pymupdf} and pdfplumber~\cite{pdfplumber}.

\begin{table}[H]
\centering
\caption{Comparison of document parsing tools}
\label{tab:docling_comparison}
\renewcommand{\arraystretch}{1.4}
\begin{tabular}{|l|c|c|c|}
\hline
\textbf{Feature} & \textbf{Docling} & \textbf{PyMuPDF} & \textbf{pdfplumber} \\
\hline
PDF support          & \checkmark & \checkmark & \checkmark \\
DOCX support         & \checkmark & \texttimes & \texttimes \\
PPTX support         & \checkmark & \texttimes & \texttimes \\
Layout analysis      & Advanced   & Basic      & Moderate   \\
Table extraction     & Strong     & Limited    & Strong     \\
OCR support          & Yes        & No         & No         \\
Metadata extraction  & Yes        & Limited    & No         \\
\hline
\end{tabular}
\end{table}

Docling performs four successive operations on each document:

\begin{enumerate}
  \item \textbf{Content extraction} reads raw text and table structures
        while preserving the logical reading order, including complex
        multi-column layouts.
  \item \textbf{Cleaning and normalisation} removes formatting artefacts
        such as repeated headers, footers, page numbers, and
        table-of-contents entries.
  \item \textbf{Metadata extraction} captures the source file name, page
        numbers, section titles, detected language, and creation date.
        This metadata travels with every chunk and is used for citation
        generation in the final answers.
  \item \textbf{Chunking} splits each document into semantically coherent chunks using a hybrid
strategy that preserves document structure and applies overlap windows
to reduce context loss between chunks.
\end{enumerate}

% ─────────────────────────────────────────────────────────────────────────────
\subsection{Embedding Generation}
% ─────────────────────────────────────────────────────────────────────────────

Each chunk is transformed into a dense numerical vector that encodes its
semantic content in a continuous high-dimensional space, enabling
similarity-based retrieval even without exact keyword overlap.

\textbf{BGE-M3}~\cite{bge_m3} was selected as the embedding model based on
its superior multilingual performance on the MIRACL and MKQA benchmarks, its
coverage of 100+ languages (including French, English, and Arabic), and its
hybrid retrieval framework. The three retrieval modes it offers are all
leveraged in Ask-Jade:

\begin{itemize}
  \item \textbf{Dense embeddings} primary representation capturing global
        semantic meaning; best balance of efficiency and expressiveness.
  \item \textbf{Sparse embeddings} complement dense retrieval for
        keyword-sensitive queries through weighted lexical signals.
  \item \textbf{Multi-vector embeddings} applied for fine-grained
        token-level matching between query and document representations.
\end{itemize}

% ─────────────────────────────────────────────────────────────────────────────
\subsection{Vector Storage}
% ─────────────────────────────────────────────────────────────────────────────

Chunk embeddings are stored in a PostgreSQL database extended with the
pgvector extension. PostgreSQL was retained because it is already
part of the existing infrastructure, avoiding the operational overhead of an
additional external service. Efficient similarity search is enabled by an
Approximate Nearest Neighbor (ANN) index based on HNSW graphs, which quickly
surfaces chunks that are semantically close to a query even in large-scale
collections. Retrieval is further refined by metadata filters on document
type, language, and project scope.

% ═══════════════════════════════════════════════════════════════════════════════
\section{Retrieval Pipeline}
\label{sec:retrieval}
% ═══════════════════════════════════════════════════════════════════════════════

The retrieval pipeline processes a user query through six sequential stages and
returns a grounded, cited response. Figure~\ref{fig:retrieval_pipeline}
illustrates the full flow.

%\begin{figure}[H]
  %\centering
  %\includegraphics[width=\textwidth]{figure_retrieval.jpg}
  %\caption{Retrieval pipeline of the Ask-Jade system}
  %\label{fig:retrieval_pipeline}
%\end{figure}

% ─────────────────────────────────────────────────────────────────────────────
\subsection{Query Rewriting (Condense Step)}
% ─────────────────────────────────────────────────────────────────────────────
In multi-turn conversations, the Condense Step reformulates the user's
question into a self-contained query using the conversation history.
For single-turn interactions, the query is forwarded unchanged.

% ─────────────────────────────────────────────────────────────────────────────
\subsection{Multi-Query Generation}
% ─────────────────────────────────────────────────────────────────────────────

To improve semantic coverage, the reformulated query is expanded into 2 to 3
variants by an LLM. Because the same information can be expressed in multiple
ways across source documents, each variant is independently embedded and used
to query the vector store, broadening the candidate set beyond what a single
formulation would return.

% ─────────────────────────────────────────────────────────────────────────────
\subsection{Semantic Retrieval}
% ─────────────────────────────────────────────────────────────────────────────

Each query variant is encoded with BGE-M3 the same model used during
ingestion ensuring alignment in the shared semantic space. Similarity is
measured by cosine distance in PgVector, and the top-10 most similar chunks
are retrieved per variant, yielding 20 to 30 candidates in total. Dense
retrieval handles both conceptual queries and paraphrased questions equally
well, and naturally supports cross-lingual search (e.g.\ a French query
retrieving English source material).

% ─────────────────────────────────────────────────────────────────────────────
\subsection{RRF Fusion}
% ─────────────────────────────────────────────────────────────────────────────

The $N$ ranked lists are merged using Reciprocal Rank Fusion (RRF),
which assigns each document a score inversely proportional to its rank:

\[
  \mathrm{RRF\_score}(d) = \sum_{i=1}^{N} \frac{1}{k + \mathrm{rank}_i(d)}
\]

where $k = 60$ is a smoothing constant. RRF is robust to score-scale
differences across variants and naturally promotes documents that rank highly
across multiple formulations, producing a single merged candidate list.

% ─────────────────────────────────────────────────────────────────────────────
\subsection{Reranking}
% ─────────────────────────────────────────────────────────────────────────────

The fused list is passed to BGE-reranker-v2-m3, a cross-encoder that
jointly processes the query and each candidate chunk. Unlike the bi-encoder
used at retrieval time, the cross-encoder achieves finer relevance scoring and
reduces noisy results. Only the top-5 chunks are forwarded to answer
generation.

% ─────────────────────────────────────────────────────────────────────────────
\subsection{Answer Generation}
\label{sec:llm_generation}
% ─────────────────────────────────────────────────────────────────────────────

The final stage generates a grounded response from the top-5 reranked chunks
and the user query. A key design requirement is that the Administrator can
switch the generation model directly from the application interface no
code change or redeployment is needed allowing the organisation to adopt
new releases, control costs, or adjust to evolving performance benchmarks.

Three API-based models are supported (Table~\ref{tab:llm_comparison}).

\begin{table}[H]
\centering
\caption{LLMs available for generation in Ask-Jade}
\label{tab:llm_comparison}
\renewcommand{\arraystretch}{1.4}
\begin{tabular}{|l|c|c|c|c|}
\hline
\textbf{Model} & \textbf{Context Window} & \textbf{Reasoning}
               & \textbf{Citation} & \textbf{Multilingual} \\
\hline
GPT-4o            & High (128k)      & Very strong & Medium      & Strong \\
Claude Sonnet 4.0 & Very high (200k) & Strong      & Very strong & Strong \\
Llama 3 (via API) & Medium (8k)      & Moderate    & Low         & Moderate \\
\hline
\end{tabular}
\end{table}

\textbf{GPT-4o} is the default because it provides the best balance between
reasoning capability, multilingual understanding, and large-context processing
for professional document retrieval in the PPP consulting domain.

All supported models share a common provider abstraction layer: they receive
the same structured prompt and return a response in the same format. Switching
the active model therefore requires only updating a single configuration value
in the database. The prompt contains three elements:

\begin{enumerate}
  \item \textbf{System instruction} defines the assistant's role, output
        format, response language, and citation rules.
  \item \textbf{Retrieved context} the top-5 reranked chunks, each
        annotated with its source document, page number, and section title.
  \item \textbf{User query} the original question as submitted.
\end{enumerate}

The model is instructed to answer exclusively from the provided context, cite
every factual claim with its source chunk, and respond in the same language as
the user's query (French, English, or Arabic). The result is a grounded
response with traceable citations linked to the original document chunks.

% ─────────────────────────────────────────────────────────────────────────────
\subsection{Observability with Langfuse}
% ─────────────────────────────────────────────────────────────────────────────

After each answer is generated, the complete prompt trace system
instruction, retrieved context, user query, and model response is saved to
Langfuse ~\cite{langfuse}. This provides full observability over the
generation pipeline, enabling the team to inspect LLM calls, monitor token
usage and costs, evaluate answer quality, and debug unexpected outputs without
modifying any application code.



\section{Conclusion}
This chapter presented the two core pipelines of the Ask-Jade system.
The ingestion pipeline handles document acquisition, parsing, chunking,
and vector storage, while the retrieval pipeline processes user queries
through rewriting, multi-query expansion, semantic retrieval, fusion,
reranking, and answer generation. Together, these components form a
modular and scalable RAG architecture. The following chapter will present
the technical environment, the system interfaces, and the evaluation
of the final architecture.